%%%%%%%%%%%%%%%%%%%%%%%%%%%%%%%%%%%%%%%%%%%%%%%%%%%%%%%%%%%%%%%%%%%%%%%%%%
% Spectrum (JetBrains incubator) -- Founding ML Engineer. Tailored resume.
% Framing: LLM systems, agentic knowledge retrieval/RAG, developer tooling / code AI, MLOps.
% Fonts: Source Sans Pro. Accent: teal + purple. Compile: pdflatex (2 passes).
%%%%%%%%%%%%%%%%%%%%%%%%%%%%%%%%%%%%%%%%%%%%%%%%%%%%%%%%%%%%%%%%%%%%%%%%%%

\documentclass[10pt,a4paper]{article}

\usepackage[T1]{fontenc}
\usepackage[utf8]{inputenc}
\usepackage[default]{sourcesanspro}
\usepackage[a4paper,top=1.1cm,bottom=1.0cm,left=1.5cm,right=1.5cm]{geometry}
\usepackage{titlesec}
\usepackage{enumitem}
\usepackage{xcolor}
\usepackage{hyperref}
\usepackage{microtype}

% ---- Colours (teal primary, purple accent) ----
\definecolor{accent}{HTML}{117A8B}
\definecolor{company}{HTML}{6C2BB3}
\definecolor{muted}{HTML}{555555}

\hypersetup{colorlinks=true, urlcolor=accent, linkcolor=accent}

\setlength{\parindent}{0pt}
\linespread{1.0}
\pagestyle{empty}

\titleformat{\section}
  {\normalfont\large\bfseries\color{accent}\scshape}
  {}{0em}{}[{\color{accent!35}\titlerule[1pt]}]
\titlespacing*{\section}{0pt}{7pt}{4pt}

\newlist{tight}{itemize}{1}
\setlist[tight]{leftmargin=1.3em,itemsep=1pt,topsep=1pt,parsep=0pt,label=\textcolor{accent}{\small$\bullet$}}

\newcommand{\entry}[4]{%
  {\color{company}\bfseries #1}\hfill {\color{muted}#2}\par
  \vspace{1pt}
  {\itshape #3}\hfill {\color{muted}\itshape #4}\par
  \vspace{3pt}%
}

\newcommand{\sep}{\,\textcolor{accent!60}{\textbar}\,}

\begin{document}

%======================= HEADER =======================
\begin{center}
  {\fontsize{24}{26}\selectfont\bfseries Abhijeet Lokhande}\\[3pt]
  {\color{accent}\large ML Engineer: LLM Systems, Agentic Knowledge Retrieval \& Developer Tooling}\\[5pt]
  {\small
    London, United Kingdom \sep
    \href{mailto:abhijeet_lokhande@proton.me}{abhijeet\_lokhande@proton.me} \sep
    07480801382 \\[2pt]
    \href{https://www.linkedin.com/in/ablds/}{LinkedIn} \sep
    \href{https://github.com/abhijeetscode}{GitHub} \sep
    \href{https://huggingface.co/ab-ai}{Hugging Face}
  }
\end{center}
\vspace{2pt}

%======================= SUMMARY =======================
\section{Summary}
ML engineer who owns the LLM stack end to end and has shaped an AI product from the early stage. I design LLM pipelines and agents for knowledge extraction, retrieval, and reasoning over enterprise code, docs, and issue data, with prompt engineering, model selection, fine-tuning, and evals, and I deploy them at scale on the Anthropic API. I go deep on fundamentals, including a GPT-style transformer written from scratch, and enjoy zero-to-one product building. Strong Python, some Kotlin. MSc Data Science (King's College London); 67K+ HuggingFace downloads across open models I have fine-tuned and released.

%======================= COMPETENCIES =======================
\section{Core Strengths}
\begin{tight}
  \item[\textcolor{accent}{\small$\bullet$}] \textbf{LLM Systems:} model architectures, fine-tuning (LoRA), prompt engineering, LLM pipeline design \& evaluation, model selection, inference optimisation
  \item[\textcolor{accent}{\small$\bullet$}] \textbf{Agentic RAG \& Knowledge Retrieval:} knowledge extraction, context-retrieval agents, vector databases, hybrid retrieval, reasoning over code, docs \& issues
  \item[\textcolor{accent}{\small$\bullet$}] \textbf{Developer Tooling \& Code AI:} spec-to-code agents, code generation, from-scratch transformers; Python, Rust, some Kotlin
  \item[\textcolor{accent}{\small$\bullet$}] \textbf{MLOps \& Zero-to-One:} evals, observability (OpenTelemetry, Arize Phoenix), experiment tracking; early-stage AI stack ownership
\end{tight}

%======================= EXPERIENCE =======================
\section{Experience}
\entry{Built AI (FinTech \& AI Company)}{April 2022 -- Present}{AI/ML Engineer (Early Stage Hire)}{London, UK}
\begin{tight}
  \item Shaped and owned the AI/ML stack as an early-stage hire, from architecture and model selection through to production for 5+ enterprise clients.
  \item Built agents for knowledge extraction and context retrieval over enterprise customer and issue data, with RAG (hybrid lexical + semantic search), tool use, and grounding.
  \item Designed LLM pipelines with prompt engineering, model selection, and evals, and deployed LLM apps at scale on the Anthropic API, including an MCP server exposing 180+ tools.
  \item Established MLOps practices: evals, Quality of Service metrics, observability (OpenTelemetry to Arize Phoenix), and incident-triggering monitoring.
  \item Optimised inference for latency and cost (local embeddings, p95 around 175ms) and improved a modelling engine's performance by 85\% through a graph-based rearchitecture.
\end{tight}

\entry{King's College London}{March 2021 -- January 2022}{Research Associate}{London, UK}
\begin{tight}
  \item Developed a deep learning model using Graph Convolution Networks and GANs for drug discovery (graph-based reasoning over molecular structures), achieving a 63\% improvement in the novelty score of generated structures.
\end{tight}

%======================= PROJECTS =======================
\section{Open-Source \& Projects}
\begin{tight}
  \item \textbf{Spec-to-Code Agent:} autonomous ReAct-based coding agent (LangGraph + Anthropic API) that keeps specification and code aligned, generating installable, tested packages from plain-language specs; 9/9 pass rate across a polyglot benchmark (Python, Go, Rust, Java, JS). \href{https://github.com/abhijeetscode/Spec-to-Code-Agent}{View on GitHub}
  \item \textbf{Agentic Enterprise Assistant:} agent that extracts and retrieves knowledge from grounded customer and issue data, with RBAC at an MCP server, a grounding-validation gate, and auditable traces; hybrid BM25 + semantic search over Elasticsearch; OpenTelemetry to Arize Phoenix. Passed 13/13 evaluation cases. \href{https://github.com/abhijeetscode/Agentic-Enterprise-Assistant}{View on GitHub}
  \item \textbf{Transformer from Scratch in Rust:} GPT-style decoder-only transformer built from first principles (multi-head causal self-attention, pre-norm residual blocks, full training loop with SafeTensors checkpointing) using the Candle framework with Metal GPU acceleration. \href{https://github.com/abhijeetscode/transformer-rust}{View on GitHub}
\end{tight}

%======================= EDUCATION =======================
\section{Education}
\entry{King's College London}{}{MSc in Data Science}{London, UK}
\vspace{-3pt}
\entry{C-DAC, India}{}{Post Graduate Diploma in Advanced Computing (PG-DAC)}{Pune, India}
\vspace{-3pt}
\entry{University of Pune}{}{B.E. in Computer Science}{Pune, India}

\end{document}
